\section{Main}
\label{sec:main}

Neurons encode information by selectively representing different stimuli, behaviors, or cognitive states through their activity, forming the basis of the neural code \cite{DayanAbbott2001}. At the population level, these signals represent high-dimensional variables that guide behavior and cognition \cite{Averbeck2006}. Yet in freely behaving animals, a central challenge is not detecting neural–behavior relationships per se, but attributing a neuron’s activity to the correct explanatory variable when behavior is high-dimensional, correlated, and temporally structured. 

Neuronal selectivity is dynamically shaped by experience. In the hippocampus, place cells can rapidly form spatial receptive fields \cite{Sotskov2022}, and in many settings these maps can remain stable over weeks \cite{Wilson1993, Thompson1990}. At the same time, mixed selectivity is common  \cite{TyeMiller2024}, such that a single neuron reflects combinations of multiple stimuli, behaviors, or cognitive states rather than a single feature. Prefrontal neurons, for example, develop mixed selectivity for task-relevant variables through learning \cite{mix3, Mante2013}. Even canonical examples of functionally selective neurons can encode multiple variables: place cells often reflect not only location, but also time, rewards, and trajectories \cite{Zheng2021, Lian2022}. This principle extends to sensory processing: neurons in the primary visual cortex and retrosplenial cortex can encode both remembered and currently perceived images \cite{Makino2015}. Together, these findings imply that apparent selectivity must be interpreted in the context of other correlated variables that co-occur during natural behavior. 

This attribution problem becomes especially acute in naturalistic settings because behavioral variables are intrinsically correlated.  Advances in automated behavior analysis, including unsupervised (MoSeq) \cite{Weinreb2024}, supervised (DeepAction \cite{Harris2023deepaction}, MARS \cite{Segalin2021}), and rule-based (BehaviorDEPOT) \cite{Gabriel2022} approaches, now enable high-throughput extraction of dozens to hundreds of behavioral features from video recordings, producing rich yet highly correlated datasets. This creates an identifiability problem: apparent selectivity for one variable can be explained by selectivity for a correlated covariate . For instance, a neuron's activity might appear to encode locomotion when it is actually tuned to speed, just as a place cell could be misinterpreted as responding to social interaction simply because encounters are tied to specific locations. Thus, the prevalence of mixed selectivity poses a fundamental problem of interpretation. This complication is prevalent in neuroscience: uninstructed movements explain substantial cortex-wide activity \cite{Musall2019}, theta modulation confounds speed versus phase coding in the hippocampus \cite{Vanderwolf1969, McNaughton1983}. These examples underscore a general requirement for selectivity analyses: they must explicitly address correlated covariates and the temporal structure of neural and behavioral signals. 

Calcium imaging further intensifies these demands by providing large-scale recordings from hundreds to thousands of neurons during continuous behavior  \cite{Chen2013, Zong2022}. Unlike spikes, calcium fluorescence is a continuous signal with complex temporal dynamics. Even the fastest indicators (e.g., jGCaMP8) exhibit decay times of tens to hundreds of milliseconds \cite{Zhang2023}. Indicator kinetics cause temporal convolution, where fluorescent signal from earlier events spuriously correlates with later behaviors \cite{Wei2020}. Moreover, behavior is described by variables of fundamentally different types---from discrete categorical states like grooming or rearing to continuous measures like speed or head direction---requiring analytical approaches that can handle heterogeneous variable types within a single framework. Many spike-based analyses assume point-process observations and do not transfer directly to temporally filtered fluorescence signals, yet deconvolution can introduce model-dependent distortions and may reduce performance in downstream decoding or attribution analyses  \cite{Rupprecht2021}.

Many existing approaches address only subsets of these requirements. Linear measures fail to detect nonlinear dependencies, dimensionality reduction obscures individual neuronal contributions, and probabilistic decoders may predict complex behavior without quantifying which specific behavioral variables explain a given neuron’s activity—especially when covariates are correlated  \cite{Kriegeskorte2019}.
Information-theoretic approaches offer a principled language for quantifying neural–behavior relationships in continuous data, but current toolboxes largely target different experimental regimes or analysis goals. The most comprehensive, the Neuroscience Information Toolbox (NIT) \cite{Maffulli2022}, provides multiple MI estimators with bias correction and detailed guidance for calcium imaging, but is designed for trial-based experimental paradigms with many repeated trials per condition. By contrast, continuous free-behavior recordings require preserving temporal autocorrelation, estimating optimal neuron–behavior delays, and controlling for covariance among behavioral features. Other toolboxes emphasize complementary levels of analysis : MINT for population-level information transmission \cite{MINT2025}; FRITES for group-level inference in electrophysiological data \cite{Frites2022}; IDTxl for network connectivity \cite{IDTxl2019}; and HOI for higher-order interactions analysis \cite{Neri2024}. The Gaussian Copula Mutual Information (GCMI) framework enables efficient MI estimation for continuous data \cite{gcmi}, but MI estimation alone does not resolve key design challenges of free behavior - lag selection under autocorrelation and disambiguation of encoding from correlation-driven artifacts. Consequently, a gap remains: no existing tool provides information-theoretic selectivity analysis that jointly handles continuous calcium fluorescence, mixed behavioral variable types, temporal delay optimization, and disentanglement of true encoding from correlation artifacts---at the scale of thousands of neurons and dozens of behavioral features.

Here we present INTENSE (INformation-Theoretic Evaluation of Neuronal SElectivity), an open-source tool that integrates efficient MI estimation with analysis components tailored to continuous free-behavior calcium imaging. INTENSE enables information-theoretic assessment of single-neuron selectivity at scale while accounting for temporal delays, heterogeneous behavioral variables, and correlated behavioral feature spaces. 

